% Template last modified by Jake Hart; please contact course staff if you have any questions regarding using this template

\documentclass{cisXXX} % You must have the cisXXX .cls file in your project or working directory (i.e. the same directory as this document) 

\HWauthor{Simone Rinaldi 737270}{s.rinaldi012@studenti.unibs.it} % Put your name and Penn email on this line
% \HWno{} % Enter the number of the homework you are working on
\HWcourse{Algoritmi e Strutture Dati} % Enter the course department and number here
\date{Elaborato a.a. 2025/26 \\[0.3em] Prof.ssa Marina Zanella}
% \HWpartner{} % If your class allows group work, put your partners here
% \HWpartner{Amelia Earhart} % Otherwise, delete or comment these lines 

\usepackage{amsmath}
\usepackage{amssymb}                                                                                                                    
\usepackage[italian]{babel}

\begin{document}
\maketitle

\section*{Compito 1}                                                                           
\subsection{Descrizione del problema}                                                                                                                                                                               
Il problema affrontato in questo elaborato consiste nell'ottimizzazione dell'ingaggio di giocatori per una squadra sportiva a budget      
limitato. Tale formulazione modella il tipico scenario decisionale affrontato dalla dirigenza di un club durante una sessione di            
calciomercato: la necessità di comporre la rosa più competitiva possibile rispettando vincoli di bilancio e di copertura dei ruoli.         
																																			
In letteratura accademica, il problema è tratto dal testo di riferimento \textit{Introduction to Algorithms} (T.~H.~Cormen, C.~E.         
~Leiserson, R.~L.~Rivest, C.~Stein), in particolare dal Problema 15-12 (``Signing free-agent baseball players'') presente nel capitolo      
dedicato alla programmazione dinamica. Nella teoria dell'ottimizzazione combinatoria e della ricerca operativa, tale problema costituisce   
un'istanza canonica del Multiple-Choice Knapsack Problem (MCKP), variante specializzata del classico problema dello zaino (0/1 Knapsack) in 
cui gli oggetti sono partizionati in classi disgiunte e mutualmente esclusive.

\subsection{Formulazione Matematica dell'Input}                                                                                           
                                                                                                                                              
    L'istanza del problema riceve in ingresso un insieme strutturato di dati discreti e deterministici:                                       
                                                                                                                                              
    \begin{itemize}                                                                                                                           
        \item \textbf{Budget Totale ($B$):} un intero non negativo $B \in \mathbb{N}$ che rappresenta la disponibilità finanziaria massima    
  spendibile per l'ingaggio complessivo dei componenti della squadra.                                                                         
        \item \textbf{Numero di Ruoli ($N$):} un intero positivo $N \in \mathbb{N}$ che definisce il numero complessivo di ruoli/reparti      
  tecnici necessari alla formazione della squadra (es. portiere, difensore, centrocampista, attaccante). I ruoli sono indicizzati da $i = 1, 2,
  \dots, N$.                                                                                                                                  
        \item \textbf{Candidati per Ruolo ($P_i$):} per ciascun ruolo $i \in \{1, \dots, N\}$, è disponibile sul mercato un insieme finito di 
  $P_i$ candidati svincolati, indicizzati con $k \in \{1, 2, \dots, P_i\}$.                                                                   
        \item \textbf{Parametri dei Giocatori:} ogni candidato $k$ appartenente al ruolo $i$ è univocamente caratterizzato da una coppia di   
  parametri non negativi:                                                                                                                     
        \begin{itemize}                                                                                                                       
            \item $c_{i,k} \in \mathbb{N}$: costo d'acquisto/ingaggio del giocatore.                                                          
            \item $v_{i,k} \in \mathbb{R}^+$ (o $\mathbb{N}$): valore tecnico atteso o metrica di rendimento del giocatore (corrispondente al 
  parametro VORP -- \textit{Value Over Replacement Player} del testo di Cormen).                                                              
        \end{itemize}                                                                                                                         
    \end{itemize}                                                                                                                             
                                                                                                                                              
                                                                                                                                              
    \subsection{Vincoli del Problema}                                                                                                         
                                                                                                                                              
    Una selezione di giocatori è considerata ammissibile se e solo se soddisfa simultaneamente i seguenti due vincoli operativi:              
                                                                                                                                              
    \begin{itemize}                                                                                                                           
        \item \textbf{Vincolo di Rosa (Esclusività di Ruolo):} per ogni singolo ruolo $i \in \{1, \dots, N\}$, è ammesso l'ingaggio di al     
  massimo un giocatore. È consentito non ingaggiare alcun candidato in un determinato ruolo (mantenendo la riserva d'ufficio presente in rosa 
  a costo $0$ e apporto incrementale nullo).                                                                                                  
                                                                                                                                              
        Formalmente, introducendo le variabili binarie di decisione:                                                                          
        \[                                                                                                                                    
        x_{i,k} \in \{0, 1\} \quad \forall i \in \{1, \dots, N\}, \quad \forall k \in \{1, \dots, P_i\}                                       
        \]                                                                                                                                    
        il vincolo si esprime come:                                                                                                           
        \[                                                                                                                                    
        \sum_{k=1}^{P_i} x_{i,k} \le 1 \quad \forall i \in \{1, \dots, N\}                                                                    
        \]                                                                                                                                    
                                                                                                                                              
        \item \textbf{Vincolo di Bilancio (Budget Globale):} la somma cumulativa dei costi associati a tutti i giocatori ingaggiati non deve  
  eccedere la soglia di budget $B$ fissata a monte:                                                                                           
        \[                                                                                                                                    
        \sum_{i=1}^N \sum_{k=1}^{P_i} c_{i,k} \cdot x_{i,k} \le B                                                                             
        \]                                                                                                                                    
    \end{itemize}                                                                                                                             
                                                                                                                                              
                                                                                                                                              
    \subsection{Funzione Obiettivo e Output Richiesto}                                                                                        
                                                                                                                                              
    Il problema si qualifica come un problema di massimizzazione vincolata.                                                                   
                                                                                                                                              
    \begin{itemize}                                                                                                                           
        \item \textbf{Funzione Obiettivo:} determinare l'assegnamento ottimale delle variabili decisionali $x_{i,k}$ tale da massimizzare il  
  valore cumulativo (competitività totale della rosa) apportato alla squadra:                                                                 
        \[                                                                                                                                    
        \max \sum_{i=1}^N \sum_{k=1}^{P_i} v_{i,k} \cdot x_{i,k}                                                                              
        \]                                                                                                                                    
        sotto il rispetto dei vincoli definiti alla sottosezione 1.3.                                                                         
                                                                                                                                              
        \item \textbf{Output Richiesto:} conformemente alle specifiche di progetto, la risoluzione deve produrre:                             
        \begin{enumerate}                                                                                                                     
            \item \textbf{Il valore scalare ottimo}, ovvero il massimo VORP complessivo conseguibile senza violare il vincolo finanziario.    
            \item \textbf{La configurazione ottima} (la rosa definitiva), descritta mediante l'elenco esatto dei giocatori scelti ruolo per   
  ruolo che realizzano tale rendimento.                                                                                                       
        \end{enumerate}                                                                                                                       
    \end{itemize}



\HWproblem
This is the first problem.
This is the first problem.
This is the first problem.
This is the first problem.
This is the first problem.
This is the first problem.

This is the second paragraph of the first problem. This equation $c^2 = a^2 + b^2 - 2ab \cos(\theta_C)$ is an example of in-line math.

And this is an example of display mode math:
$$x = \frac{-b \pm \sqrt{b^2 - 4ac}}{2a}$$
This is an example of the align* environment, which can be helpful for mathematical expressions with multiple steps:
\begin{align*}
r &= 1/2 + 1/3 + 1/4 + 1/5 + 1/6 + 1/7 + 1/8 + \ldots + 1/n\\
&> 1/2 + 1/4 + 1/4 + 1/8 + 1/8 + 1/8 + 1/8 + \ldots + 1/n\\
&= \sum_{i = 1}^{\log_2 n} \frac{2^{i - 1}}{2^i}\\
&= \frac{1}{2} \log_2 (n)\\ 
&= \Omega(\log n)
\end{align*}
This is an example of matrices in \LaTeX:
\begin{align*}
	\sigma_x &=
		\begin{pmatrix}
		0 & 1 \\
		1 & 0
		\end{pmatrix} \\
	\sigma_y &=
		\begin{pmatrix}
		0 & -i \\
		i & 0
		\end{pmatrix} \\
	\sigma_z &=
		\begin{pmatrix}
		1 & 0 \\
		0 & -1
		\end{pmatrix} \\
\end{align*}
This concludes the first problem and the mathematical examples. The following questions demonstrate how to use the template for multipart problems.

\HWproblem
\HWsubproblem
This is the second problem, first question.
This is the second problem, first question.
This is the second problem, first question.
This is the second problem, first question.
This is the second problem, first question.
This is the second problem, first question.
\HWsubproblem
This is the second problem, second question.
This is the second problem, second question.
This is the second problem, second question.
This is the second problem, second question.
This is the second problem, second question.
This is the second problem, second question.
\HWsubproblem
This is the second problem, third question.
This is the second problem, third question.
This is the second problem, third question.
This is the second problem, third question.
This is the second problem, third question.
This is the second problem, third question.


\HWproblem
This is an introduction to the the third problem.
\HWsubproblem
This is the third problem, first question.
This is the third problem, first question.
This is the third problem, first question.
This is the third problem, first question.
This is the third problem, first question.
This is the third problem, first question.
\HWsubproblem
This is the third problem, second question.
This is the third problem, second question.
This is the third problem, second question.
This is the third problem, second question.
This is the third problem, second question.
This is the third problem, second question.
\HWsubproblem
This is the third problem, third question.
This is the third problem, third question.
This is the third problem, third question.
This is the third problem, third question.
This is the third problem, third question.
This is the third problem, third question.
\HWsubproblem
This is the third problem, fourth question.
This is the third problem, fourth question.
This is the third problem, fourth question.
This is the third problem, fourth question.
This is the third problem, fourth question.
This is the third problem, fourth question.

\HWproblem
And this is the final one.
And this is the final one.
And this is the final one.
And this is the final one.
And this is the final one.
And this is the final one.

And this is the final one.
And this is the final one.
And this is the final one.
And this is the final one.
And this is the final one.
And this is the final one.

And this is the final one.
And this is the final one.
And this is the final one.
And this is the final one.
And this is the final one.
And this is the final one.
\end{document}
